\section{Background}
\label{sec:background}

\paragraph{Derivation tree.} A grammar derives a query by choosing one
alternative per rule. The \emph{derivation tree} \Tree{\Q} of a query \Q
records these choices: each inner node is a rule with the alternative it took,
and the leaves are SQL tokens. Reading the leaves from left to right gives the
text of \Q.

\paragraph{Wrong result.} A query has a wrong result when the server returns a
multiset of rows \Res{\Q} that differs from what the SQL standard and the \pg
documentation define, without an error. Typical causes are planner rules that
are correct only under an assumption that does not hold, for example that a
column cannot be NULL.

\paragraph{Metamorphic oracle.} A metamorphic oracle builds a second query
$\Q'$ with a known relation to \Q, usually $\Res{\Q} = \Res{\Q'}$, and reports a
bug when the relation fails. TLP~\cite{rigger2020tlp} uses
$\Res{\Q} = \Res{\Part{\Q}{p}} \uplus \Res{\Part{\Q}{\neg p}} \uplus
\Res{\Part{\Q}{p \;\mathtt{IS\;NULL}}}$ for a predicate $p$ added to
\texttt{WHERE}. NoREC~\cite{rigger2020norec} compares the row count of
\texttt{SELECT * FROM t WHERE p} with the number of rows on which $p$ is true
when $p$ is evaluated in the select list, where the planner cannot use it to
choose a plan.

\paragraph{Plan variant.} \pg lets a session switch planner methods off, for
example \texttt{SET enable\_hashjoin = off}. A \emph{plan variant} is the
same query under such a setting. All plan variants of a deterministic query
must return the same multiset. This is differential testing inside one
server.

\paragraph{Assertion build.} \pg built with \texttt{--enable-cassert} checks
internal invariants with \texttt{Assert}. A failed assertion aborts the backend
and prints a \texttt{TRAP} line with the file and the expression to the
server log. Release builds skip these checks, so an assertion failure often
corresponds to silent corruption or a wrong result in production.
